% ============================================================
%  LETTRE DE FIN DE MISSION ESN — Phase 5 : Clôture
%  Time'Eats — Cellule PMO
% ============================================================
\documentclass[11pt,a4paper,oneside]{report}
\usepackage{../style/consulting}
\hypersetup{pdftitle={Lettre de fin de mission ESN — Time'Eats PMO}}
\pagestyle{fancy}\fancyhf{}
\renewcommand{\headrulewidth}{0pt}
\fancyhead[L]{\begin{tikzpicture}[remember picture,overlay]
  \draw[cnNavy,line width=0.5pt](0,-0.42cm)--(\textwidth,-0.42cm);
\end{tikzpicture}\vspace{2pt}\timeeatslogo\hspace{4pt}\small\textcolor{cnSlate}{Time'Eats \textbf{·} Lettre de fin de mission ESN — \textit{[Nom projet]}}}
\fancyfoot[C]{\tiny\textcolor{cnSlate}{Document confidentiel — Cellule PMO — CESI MAALSI 2026}}
\fancyfoot[R]{\small\textcolor{cnSlate}{\thepage}}
\begin{document}\small

\begin{center}
{\LARGE\bfseries\color{cnNavy} Lettre de fin de mission --- Prestataire ESN}\\[4pt]
\textcolor{cnOrange}{\rule{10cm}{1pt}}\\[4pt]
{\large\color{cnSlate} Phase 5 --- Clôture \quad|\quad Projet : \textit{[Nom du projet]}}
\end{center}

\vspace{4pt}
\begin{decision}
Ce document \textbf{contractuel} entérine la fin de la prestation de l'ESN. Il conditionne le \textbf{solde du forfait} et la \textbf{décharge de responsabilité}. Il est cosigné par le chef de projet Time'Eats et le responsable de mission ESN, puis archivé au dossier projet et transmis aux Achats.
\end{decision}

\vspace{4pt}
\renewcommand{\arraystretch}{1.4}
\begin{tabularx}{\textwidth}{L{3.5cm} Y L{3.5cm} Y}
\toprule
\rowcolor{cnNavy}\th{Champ} & \th{Valeur} & \th{Champ} & \th{Valeur} \\
\midrule
Projet            & \textit{[Nom]}              & Prestataire (ESN) & \textit{[Raison sociale]} \\
\rowcolor{cnIce}
N° de contrat     & \textit{[Réf. contrat]}     & Type d'engagement & \textit{[Forfait / Régie]} \\
Responsable ESN   & \textit{[Prénom NOM]}       & Chef de projet    & \textit{[Prénom NOM]} \\
\rowcolor{cnIce}
Début de mission  & \textit{[JJ/MM/AAAA]}       & Fin de mission    & \textit{[JJ/MM/AAAA]} \\
\bottomrule
\end{tabularx}

\section*{1 — Périmètre réalisé et livrables acceptés}

\renewcommand{\arraystretch}{1.3}
\begin{tabularx}{\textwidth}{C{0.6cm} L{5cm} Y C{3cm}}
\toprule
\rowcolor{cnNavy}\th{\#} & \th{Livrable contractuel} & \th{PV de référence} & \th{Statut} \\
\midrule
1 & \textit{[ex. Application web v1]}       & \textit{[PV recette fonctionnelle P5-19]} & \cellcolor{cnGreen!20}Accepté \\
\rowcolor{cnIce}
2 & \textit{[ex. API REST documentée]}      & \textit{[PV recette technique P5-20]}     & \cellcolor{cnGreen!20}Accepté \\
3 & \textit{[ex. Transfert d'exploitation]} & \textit{[PV transfert P5-22]}             & \cellcolor{cnGreen!20}Accepté \\
\bottomrule
\end{tabularx}

\section*{2 — Levée des réserves}

\renewcommand{\arraystretch}{1.3}
\begin{tabularx}{\textwidth}{C{0.6cm} Y C{3cm} C{3cm}}
\toprule
\rowcolor{cnNavy}\th{\#} & \th{Réserve émise en recette} & \th{Date de levée} & \th{Statut} \\
\midrule
1 & \textit{[Description de la réserve]} & \textit{[JJ/MM/AAAA]} & \cellcolor{cnGreen!20}Levée \\
\rowcolor{cnIce}
2 & \textit{[Description de la réserve]} & \textit{[JJ/MM/AAAA]} & \cellcolor{cnGreen!20}Levée \\
\multicolumn{4}{l}{\footnotesize \textit{Aucune réserve bloquante ouverte à la date de signature.}} \\
\bottomrule
\end{tabularx}

\section*{3 — Solde financier et clôture contractuelle}

\renewcommand{\arraystretch}{1.3}
\begin{tabularx}{\textwidth}{L{6cm} C{3cm} Y}
\toprule
\rowcolor{cnNavy}\th{Élément} & \th{Montant} & \th{Commentaire} \\
\midrule
Montant du forfait contractualisé & \textit{[XXX K€]} & \\
\rowcolor{cnIce}
Déjà facturé / réglé              & \textit{[XXX K€]} & \\
Solde à facturer (dernier jalon)  & \textit{[XX K€]}  & À réception de la présente lettre signée \\
\rowcolor{cnIce}
Pénalités éventuelles             & \textit{[0 / XX K€]} & \textit{[Retard, SLA non tenu, etc.]} \\
\bottomrule
\end{tabularx}

\vspace{6pt}
\begin{reco}
\textbf{Décharge.} Sous réserve du bon paiement du solde, les parties conviennent que la mission est \textbf{achevée}. La \textbf{garantie contractuelle} (correction des anomalies) court jusqu'au \textit{[JJ/MM/AAAA]}. Les engagements de \textbf{confidentialité} et de \textbf{réversibilité} survivent à la fin de la mission.
\end{reco}

\vspace{8pt}
\renewcommand{\arraystretch}{1.6}
\begin{tabularx}{\textwidth}{L{5cm} Y Y}
\toprule
\rowcolor{cnNavy}\th{} & \th{Pour Time'Eats (CP / PMO)} & \th{Pour l'ESN (Resp. mission)} \\
\midrule
Nom \& fonction & \textit{[Prénom NOM]} & \textit{[Prénom NOM]} \\
Date           & \textit{[JJ/MM/AAAA]} & \textit{[JJ/MM/AAAA]} \\
Signature      &                       & \\
\bottomrule
\end{tabularx}

\end{document}
